% =============================================================================
% Sperrvermerk -- nur noetig, wenn deine Arbeit unternehmensinterne Daten
% enthaelt. Einbinden in main.tex vor dem Inhaltsverzeichnis.
%
% Die Leitfaeden widersprechen sich hier:
%   JKS: roemische Seitenzahl UND im Inhaltsverzeichnis
%   WI : keine Seitenzahl und NICHT im Inhaltsverzeichnis
% Der Formalia-Schalter macht automatisch das Richtige.
%
% \myFirma in skripte/meta.tex setzen.
% Uebernommen aus dem FOM-LaTeX-Template von Andy Grunwald (MIT) und angepasst.
% =============================================================================
\clearpage
\ifformaliaSperrvermerkImTOC
  \addcontentsline{toc}{section}{Sperrvermerk}
\else
  \thispagestyle{empty}
\fi

\section*{Sperrvermerk}

Die vorliegende Arbeit mit dem Titel \enquote{\myTitel} enthält
unternehmensinterne Daten der Firma \myFirma. Sie ist daher nur zur Vorlage
bei der FOM Hochschule sowie bei den Begutachtenden der Arbeit bestimmt. Für
die Öffentlichkeit und für dritte Personen darf sie nicht zugänglich sein.

\vspace{4cm}

\noindent
\begin{tabular*}{\textwidth}{@{}l@{\extracolsep{\fill}}r@{}}
  \rule[0.5ex]{12em}{0.55pt} & \rule[0.5ex]{12em}{0.55pt} \\
  (Ort, Datum) & (Eigenhändige Unterschrift) \\
\end{tabular*}
